% The Journal of Portfolio Management submission.
% Built to the PMR Article Submission Guidelines (March 2023):
%   - target 4,000 words of main text, excluding abstract, exhibits, references
%   - abstract 160 words, non-technical, no references
%   - three-item key takeaways, each under 40 words
%   - tabular material is called an Exhibit, never a Table
%   - Chicago Manual of Style author-date references
%   - Times New Roman 12pt, 1in margins, single spaced, 12pt before paragraphs
%   - page 1 front matter, page 2 abstract block, main text opens page 3
\documentclass[12pt]{article}

\usepackage[T1]{fontenc}
\usepackage[utf8]{inputenc}
\usepackage{times}                    % Times New Roman
\usepackage[margin=1in]{geometry}
\usepackage{booktabs}
\usepackage{dcolumn}   % decimal-point alignment, per PMR exhibit standards
\newcolumntype{d}[1]{D{.}{.}{#1}}
\usepackage{amsmath}
\usepackage[hidelinks]{hyperref}

\setlength{\parindent}{0.5in}         % first line indented one tab stop
\setlength{\parskip}{12pt plus 2pt}   % 12 points of space before paragraphs
\linespread{1.0}                      % single spaced

\renewcommand{\footnotesize}{\fontsize{10}{12}\selectfont}

% Headings, per PMR: top level bold uppercase, second level indented and bold.
% Defined manually rather than via titlesec, which this distribution lacks.
\makeatletter
% Top level: bold, uppercase, on its own line.
\renewcommand{\section}[1]{\par\vspace{16pt}%
  \noindent{\bfseries\MakeUppercase{#1}}\par\nobreak\vspace{4pt}%
  \@afterindenttrue\@afterheading}
% Second level: indented and bold, on its own line.
\renewcommand{\subsection}[1]{\par\vspace{10pt}%
  \noindent\hspace*{0.5in}{\bfseries #1}\par\nobreak\vspace{2pt}%
  \@afterindenttrue\@afterheading}
\makeatother
% Third level: bold run-in heading; the paragraph continues on the same line.
\newcommand{\runin}[1]{\par\vspace{2pt}\noindent\hspace*{0.5in}{\bfseries #1}\ \ignorespaces}

% Exhibits are not floats: PMR numbers them in order of appearance and the text
% must refer to each one, so fixing them in place keeps text and exhibit adjacent.
% Wrapped in a minipage so the title can never be orphaned from its table.
\newenvironment{exhibit}[2]{%
  \par\vspace{14pt}\noindent\begin{minipage}{\textwidth}\centering\small
  \textbf{Exhibit #1: #2}\par\vspace{6pt}}%
  {\end{minipage}\par\vspace{12pt}}
\newcommand{\exnote}[1]{\par\vspace{4pt}%
  \begin{minipage}{0.88\textwidth}\fontsize{9}{11}\selectfont #1\end{minipage}}

\pagestyle{plain}                     % page numbers at bottom centre

\begin{document}

% ============================ PAGE 1 ============================
\thispagestyle{empty}
\vspace*{1in}
\begin{center}
{\large\bfseries Market-Neutral Alpha in Crypto Perpetual Futures:\\
Ranking Each Asset Against Itself}

\vspace{24pt}
August 2026

\vspace{36pt}
\textbf{Dhanya MD}\\
Independent Researcher\\
% JPM's submission form requires a mailing address and telephone on page 1,
% although their published style guide mentions neither.
Jayanagar, Hassan\\
Karnataka 543201, India\\
+91 93534 96617\\
\href{mailto:dhanya13md@gmail.com}{dhanya13md@gmail.com}

\vspace{24pt}
\begin{minipage}{0.8\textwidth}\centering
Dhanya MD is an independent quantitative researcher.\\
Corresponding author.
\end{minipage}
\end{center}

\newpage

% ============================ PAGE 2 ============================
\thispagestyle{empty}
\begin{center}
{\large\bfseries Market-Neutral Alpha in Crypto Perpetual Futures:\\
Ranking Each Asset Against Itself}

\vspace{12pt}
August 2026
\end{center}

\vspace{12pt}
\noindent\textbf{ABSTRACT}

\noindent Factor investing in equities works by comparing companies to one
another. Shared accounting rules, industry membership and cash flow make those
comparisons meaningful. Cryptocurrency perpetual futures share none of that, yet
practitioners rank them against each other anyway. We test what happens when you
stop. Applying ten factors drawn from published equity research to 112 crypto
contracts over seven years, we find that scoring each contract against its own
trailing history, rather than against its peers, raises average annual Sharpe
from 1.01 to 1.45. The improvement holds for nine of the eleven factors tested
and survives removing the funding component entirely. Building one portfolio per
factor and blending returns, rather than blending scores into a single ranking,
adds a further 0.67 Sharpe. The combined strategy returned 27\% annually with a
7.6\% maximum drawdown and gained 9\% during 2022, when the market fell 82\%.

\vspace{18pt}
\noindent\textbf{KEY TAKEAWAYS}

\begin{itemize}\setlength{\itemsep}{6pt}\setlength{\parskip}{0pt}
\item Crypto contracts have no shared fundamentals, so ranking them against each
other compares numbers that do not mean the same thing. Ranking each against its
own history fixes this and adds 0.44 Sharpe.

\item Funding is not a cost adjustment in perpetual futures. It is a return
driver worth 40\% to 266\% of each factor's performance, and a price-only
backtest understates results.

\item Holding one book per factor and blending the returns beats blending signals
into a single score by 0.67 Sharpe, with both approaches tuned over the same
1,286-configuration grid.
\end{itemize}

\vspace{18pt}
\noindent\textbf{Keywords:} market neutral, cryptocurrency, perpetual futures,
factor investing, cross-sectional ranking, portfolio construction

\noindent\textbf{JEL:} G11, G12, G14, C58

\newpage

% ============================ PAGE 3 ============================
\indent Every cross-sectional strategy rests on an assumption that is rarely stated out
loud: that the assets being ranked are comparable. When a portfolio manager sorts
equities on book-to-price and buys the cheapest quintile, the sort is meaningful
because a price-to-book ratio of 0.8 means roughly the same thing at one
industrial company as at another. Shared accounting standards, industry
classification and observable cash flow supply that common ground.

Cryptocurrency perpetual futures supply none of it. There is no accounting
identity linking one token to another, no sector taxonomy with economic content,
no cash flow to normalise by. A market-capitalisation figure depends on a token
release schedule set by a founding team. Trading volume spans four orders of
magnitude across contracts on the same exchange. Yet the strategies practitioners
run in this market are almost uniformly cross-sectional: rank every contract on a
signal, buy the top, sell the bottom.

The cross-sectional approach has deep roots in equities, from Fama and French
(1993) through Asness, Moskowitz and Pedersen (2013), and it travels well across
equity markets precisely because the comparability assumption holds there. Recent
work extends factor pricing to digital assets (Cong et al. 2026), but does so by
manufacturing a comparability anchor rather than questioning whether one exists.

This article asks what that assumption costs, and what happens if you abandon it.
The alternative is straightforward. Instead of asking whether a contract's
momentum is high relative to other contracts, ask whether it is high relative to
its own history. Each asset becomes its own benchmark. The comparison no longer
requires the assets to be alike, only that each has a past. The idea has a
precedent in time-series momentum (Moskowitz, Ooi and Pedersen 2012), which
judges each instrument against its own trailing return rather than against peers;
we apply the same logic to the ranking step of a cross-sectional book.

We test this on 112 US-dollar-margined perpetual futures over 363 weeks, using
ten factors taken from published equity research plus one built specifically for
this market. The change in ranking frame is worth 0.44 in annual Sharpe ratio,
holding in nine of eleven cases. A second construction choice, holding a separate
portfolio for each factor rather than blending signals into one score, is worth a
further 0.67. Together they turn a set of individually unremarkable signals into a
strategy that returned 27\% annually with a maximum drawdown under 8\%.

The remainder of this article describes the market's peculiarities, the strategy,
what each design decision contributes, and the evidence that the result is not an
artefact of testing too many ideas against too little data.

\section{Why This Market Is Different}

Three features of perpetual futures matter for portfolio construction.

\runin{No anchor for comparison.}The absence of shared fundamentals is not merely inconvenient. It means the same
raw factor value carries different information at different contracts. A
volatility reading of 30\% annualised is unremarkable for a small-cap token and
extraordinary for the largest contract on the exchange. A cross-sectional sort
treats those two readings as identical evidence. They are not.

\runin{Funding is a return, not a fee.}Perpetual futures have no expiry. They track spot through a funding payment
exchanged between long and short holders, typically every eight hours. Most
practitioners treat this as a carrying cost to be netted out. That is a mistake
here. Over our sample the funding component contributes between 40\% and 266\% of
each factor's total return, and a book constructed as ours is receives funding in
98\% of weeks. A backtest run on price alone understates this strategy rather than
flattering it.

\runin{Aggressor-side information is published.}Centralised crypto venues publish which side of each trade was the aggressor,
information that equity researchers must infer with classification algorithms
(Lee and Ready 1991). This makes it possible to measure something equities cannot
observe directly: whether buy-initiated trades arrive in a different size
distribution than sell-initiated ones. We introduce a factor built from
the difference in skewness between the two. We are not aware of prior work that
uses the skewness of the trade-size distribution, measured separately by
initiating side and differenced over time, as a cross-sectional return predictor.
Zhou (2008) partitions the size distribution by side in an order-driven equity
market but measures tail exponents, finding them statistically indistinguishable
across sides; a separate literature prices the skewness of \emph{returns}. Neither
measures the skewness of trade \emph{sizes}, nor how the asymmetry between sides
evolves.

The construction is simple enough to state in one line. Within each five-minute
bar, classify the bar as buy-dominated or sell-dominated by aggressor volume,
take the log of the average trade size in each, compute the skewness of each
series over the week, and difference the two:
%
\begin{equation*}
\text{TSKD}_t \;=\; \Delta\Big[\,\mathrm{skew}\big(\log \text{size}^{B}\big)
 \;-\; \mathrm{skew}\big(\log \text{size}^{S}\big)\Big]
\end{equation*}
%
The differencing matters: the level of the asymmetry is a persistent property of a
contract, while the \emph{change} in it is what carries information. The measure
requires five-minute resolution and cannot be recovered from daily bars, which is
why it is absent from the equity literature that relies on inferred trade
direction.

The intuition is that informed and uninformed order flow have different size
signatures. A participant with a view accumulates in a pattern that differs from
retail flow reacting to a price move, and the difference shows up in the tail of
the size distribution rather than in its average. Averages are what order-imbalance
measures capture, and they discard exactly the information that distinguishes the
two. Measuring the skewness gap between buy-dominated and sell-dominated intervals
retains it.

\section{The Factors}

Eleven factors are used. Ten are transported from published equity research and the
formulaic-alpha literature (Kakushadze 2016)
rather than invented here, which is deliberate: the article's claim concerns how
factors are ranked and combined, not whether we can find new ones. Using
established signals means any improvement is attributable to construction rather
than to signal discovery.

They fall into four families. \emph{Volatility and jump measures} capture the
character of recent price movement, distinguishing realised upside from downside
variation (Barndorff-Nielsen, Kinnebrock and Shephard 2010) and separating
continuous moves from jumps (Patton and Sheppard 2015). \emph{Liquidity and
impact measures} relate price movement to the volume required to produce it,
including a spread proxy and a quadratic-impact term. \emph{Order-flow measures}
use the aggressor-side data described above, including signed imbalance (Chordia, Roll and
Subrahmanyam 2002) and the
trade-size asymmetry factor; the premise that order size carries information
traces to Kyle (1985). \emph{Positioning measures} use exchange-published
statistics on open interest and the long-short ratios of top traders relative to
the broader account base.

Individually, none of these is remarkable. Traded alone in the conventional
cross-sectional frame, the eleven average an annualised Sharpe of 1.01, which
after costs is not a business. The argument of this article is that most of the
value in this market sits in decisions taken after the signals are chosen.

Exhibit 1 lists them with the data each requires.

\begin{exhibit}{1}{The Eleven Factors and Their Data Requirements}
\begin{tabular}{llc}
\toprule
Family & Factor & Data required \\
\midrule
Volatility and jumps & Realised semivariance ratio & daily bars \\
 & Signed jump variation & daily bars \\
 & Adjusted volatility & daily bars \\
Liquidity and impact & Spread proxy & daily bars \\
 & Quadratic impact & daily bars \\
 & Volume-scaled impact, mean & daily bars \\
 & Volume-scaled impact, variance & daily bars \\
Order flow & Order-flow imbalance & daily bars \\
 & Trade-size skewness gap & \textbf{5-minute} \\
Positioning & Open-interest change & exchange metrics \\
 & Top-trader vs.\ all-account ratio & exchange metrics \\
\bottomrule
\end{tabular}
\exnote{Ten of the eleven are transported from published equity research. Only
the trade-size skewness gap requires intraday resolution; a practitioner working
from daily bars can implement the other ten. Exchange metrics are published
free and require no credentials.}
\end{exhibit}

A practitioner without five-minute history can implement ten of the eleven and
lose relatively little. The trade-size factor contributes positively but is not
load-bearing.

\section{The Strategy}

The construction has four steps.

\runin{Score each contract against its own past.}For each of eleven factors, a contract's raw factor value in a given week is
converted to a percentile against that contract's own trailing 52-week
distribution of the same factor. The score is centred so that zero is the
contract's own median. Nothing in the calculation refers to any other contract.
Critically, the trailing window uses only weeks strictly before the scoring week,
so a score can be computed in real time with no knowledge of the future.

To see the difference this makes, consider two contracts in a single week of our
sample. One shows a raw factor value of $-27.86$, the lowest in the
cross-section. The other shows $-11.74$, far higher. Cross-sectional ranking would
buy the second and sell the first. Measured against their own histories, the
first is only mildly unusual and scores near the middle, while the second is
extreme for that contract and scores near the top. The self-referential frame
reverses the conclusion, and it is the frame that proves correct.

\runin{Build one portfolio per factor.}Each factor independently sorts the tradeable universe and holds the top 20\% long
and the bottom 20\% short, equally weighted, dollar neutral. With 112 contracts
that is 22 positions a side at roughly 4.5\% each. The quintile fraction is a
parameter, not a constant: it is selected on historical data alongside everything
else, and on a smaller universe the position count shrinks with it.

\runin{Combine returns, not scores.}The eleven single-factor portfolios are combined by weighting each in inverse
proportion to its own realised volatility, so a steadier factor contributes more.
Combining weakly correlated books rather than collapsing them into one is what
Grinold's (1989) fundamental law would predict: information ratio scales with the
square root of the number of independent bets, and blending scores destroys that
breadth.
This is the step most implementations skip. The common approach averages the
eleven scores into one composite and sorts on that. Averaging scores destroys the
information that distinguishes factors from one another; a contract that one
factor loves and another hates ends up indistinguishable from a contract both are
indifferent to.

The difference is easiest to see in the final weights. In a representative week of
our sample, one contract is shorted by six of the eleven factors and held by none,
and ends with a portfolio weight of $-2.3\%$. Another is shorted by four and held
by one, ending at $-1.5\%$. A third, shorted by two and held by one, ends at
$-0.7\%$. The position sizes carry information about how much the factors agree.
A blended score would have compressed all three into a single ranking and sized
them by their position in that ranking, discarding the distinction between broad
agreement and marginal disagreement.

The volatility weighting matters for a second reason. The eleven factors have
materially different return volatilities, and equal-weighting their books would let
the noisiest one dominate the combined portfolio's risk regardless of its
contribution. Weighting by inverse volatility equalises risk contribution instead
of capital, which is standard practice in multi-strategy allocation and is
unusually valuable here because the dispersion in factor volatility is wide.

\runin{Tilt toward carry and cap turnover.}Positions are tilted toward receiving funding rather than paying it, and weekly
turnover is capped. Average turnover runs at 37\% of gross exposure per week.

\runin{Universe and rebalancing.}The tradeable universe is screened on trailing dollar volume, which removes
contracts too thin to trade at size without moving the price. The screen is
applied on trailing data only, so a contract that becomes liquid enters the
universe when it does so and not retrospectively. This detail matters more than it
sounds: applying a liquidity screen using full-sample data is one of the most
common ways a crypto backtest quietly acquires hindsight, because the contracts
that survive to be liquid at the end of the sample are exactly the ones that did
well.

Rebalancing is weekly. That frequency is a compromise: the factors decay over
several weeks, so daily rebalancing adds cost without adding information, while
monthly rebalancing gives up too much of the signal. Weekly also happens to be
long enough that execution can be worked patiently rather than taken at market,
which is what makes the cost assumptions in this article realistic.

\section{What Each Choice Contributes}

Exhibit 2 isolates the ranking frame. Every other element is held fixed and only
the comparison basis changes.

\begin{exhibit}{2}{The Ranking Frame}
\begin{tabular}{l d{1.3} d{1.3} c}
\toprule
\multicolumn{1}{l}{} & \multicolumn{1}{c}{Cross-sectional}
 & \multicolumn{1}{c}{Self-referential} & Factors improved \\
\midrule
Mean Sharpe, with funding    & 1.011 & 1.451 & 9 of 11 \\
Mean Sharpe, funding removed & 0.383 & 0.648 & 8 of 11 \\
\bottomrule
\end{tabular}

\exnote{Mean annualised Sharpe ratio across eleven single-factor portfolios, 112
contracts, 363 weeks. The second row removes the funding component entirely,
showing the effect is not a carry artefact.}
\end{exhibit}

The improvement is not explained by funding: with the funding component removed
entirely, the ranking frame still adds 0.27 and still holds in eight of eleven
cases. Nor is it a statistical curiosity. Standardising each contract against its
own history before ranking cross-sectionally, an intermediate step that removes
scale differences but keeps the peer comparison, recovers 41\% of the gap. That
identifies scale heterogeneity as one mechanism and leaves the remainder to the
frame itself.

It is worth being precise about what that 41\% means, because it points at what is
actually going wrong in the conventional approach. When a cross-sectional sort
ranks contracts on a raw factor value, it is implicitly assuming the factor has the
same distribution everywhere. It does not. A liquidity measure computed on the
largest contract and on a token trading a thousandth of its volume produce numbers
on entirely different scales, so the sort ends up ordering contracts by size rather
than by the characteristic it intended to capture. Standardising each contract
first removes that contamination, and recovering 41\% of the gap tells you scale
was roughly half the problem.

The other half is harder to remove, and it is the part that motivates abandoning
the peer comparison altogether. Even after standardisation, a cross-sectional sort
still asks which contract is most extreme relative to the others this week. In a
market where the constituents genuinely differ, that question has no stable answer,
because the composition of the cross-section changes as contracts list, delist and
move in and out of the liquidity screen. Asking instead whether each contract is
extreme relative to its own past sidesteps the composition problem entirely.

Exhibit 3 isolates the construction choice, comparing one-book-per-factor against
the conventional blended-score pipeline. Both were tuned over an identical grid of
1,286 configurations, so neither is handicapped.

\begin{exhibit}{3}{Portfolio Construction}
\begin{tabular}{l d{1.3} d{1.3}}
\toprule
\multicolumn{1}{l}{} & \multicolumn{1}{c}{Best} & \multicolumn{1}{c}{Average} \\
\multicolumn{1}{l}{} & \multicolumn{1}{c}{configuration}
 & \multicolumn{1}{c}{configuration} \\
\midrule
One book per factor, self-referential & 2.571 & 1.537 \\
Blended scores, cross-sectional       & 1.971 & 0.995 \\
\midrule
Mean difference, paired ($t=5.82$)    & \multicolumn{1}{c}{0.665} &
 \multicolumn{1}{c}{---} \\
Matched cells favouring one book      & \multicolumn{1}{c}{15 of 17} &
 \multicolumn{1}{c}{---} \\
\bottomrule
\end{tabular}

\exnote{Annualised Sharpe ratio. The paired comparison holds every parameter fixed and
varies only construction, so the difference is attributable to construction alone.}
\end{exhibit}

The ordering holds on both the best configuration and the average one, which
matters: a result that appears only at the optimum is usually a result about the
optimiser.

Carry deserves separate mention. Traded on its own the funding component returns
an annualised Sharpe of 1.25 ($t=2.91$), respectable but short of the $t>3.0$ bar
and so not by itself compelling. Its value
is as a component. Removing it costs the least-affected factor about 40\% of its
performance and the most-affected more than two-thirds. Practitioners who net
funding out of their backtests as a cost are discarding a return source.

\section{Performance}

Exhibit 4 summarises the combined strategy.

\begin{exhibit}{4}{Strategy Performance, 282 Weekly Rebalances}
\begin{tabular}{lr}
\toprule
Cumulative return & $+327.6\%$ \\
Annualised return & $+27.8\%$ \\
Annualised volatility & 12.8\% \\
Sharpe ratio & \textbf{2.16} \\
$t$-statistic & \textbf{5.03} \\
Maximum drawdown & \textbf{$-7.6\%$} \\
Recovery from maximum drawdown & 16 weeks \\
Worst week & $-6.2\%$ \\
Positive weeks & 65.5\% \\
Return to maximum drawdown & 3.7 \\
Skewness & $+1.14$ \\
Weekly turnover & 37\% \\
\bottomrule
\end{tabular}
\end{exhibit}

The $t$-statistic of 5.03 is worth pausing on. Harvey, Liu and Zhu (2016) argue
that the flood of published anomalies means a newly claimed factor should clear
$t>3.0$ rather than the conventional 2.0. This result clears that raised bar by
two-thirds, on 282 independent weekly observations.

Two features deserve attention beyond the headline. The maximum drawdown of 7.6\%
recovered within sixteen weeks, giving a return-to-drawdown ratio of 3.7. And
returns are positively skewed at $+1.14$. Most systematic strategies that produce
high Sharpe ratios do so by selling tails and exhibit negative skew; this one does
the opposite, which matters for anyone sizing the position.

The strategy's usefulness rests on a claim of market neutrality. A regression of
weekly returns on market, size and momentum factors built from the same universe
is reported in Exhibit 5. The market loading is $-0.001$ with a $t$-statistic of
$-0.12$: the book is market-neutral \emph{empirically}, not merely by
construction. Size and momentum load negatively but explain little, and the
regression $R^2$ is 0.074. The intercept is an annualised alpha of 27.05\% with a
$t$-statistic of 4.42 after Newey--West (1987) correction, accounting for 97\% of
the strategy's mean return. Three factors that explain much of what happens in
this market explain almost none of this.

\begin{exhibit}{5}{Is It Alpha? Regression on Market, Size and Momentum}
\begin{tabular}{l d{1.5} d{2.2} d{2.2}}
\toprule
\multicolumn{1}{l}{} & \multicolumn{1}{c}{Coefficient}
 & \multicolumn{1}{c}{$t$ (OLS)} & \multicolumn{1}{c}{$t$ (NW)} \\
\midrule
Alpha    &  0.00520 &  5.03 &  4.42 \\
Market   & -0.00128 & -0.13 & -0.12 \\
Size     & -0.09260 & -3.83 & -2.69 \\
Momentum & -0.06819 & -3.39 & -2.23 \\
\midrule
\multicolumn{1}{l}{Annualised alpha} & \multicolumn{3}{r}{\textbf{+27.05\%}} \\
\multicolumn{1}{l}{$R^{2}$}          & \multicolumn{3}{r}{0.074} \\
\multicolumn{1}{l}{Alpha share of mean return} & \multicolumn{3}{r}{\textbf{97\%}} \\
\multicolumn{1}{l}{Observations}     & \multicolumn{3}{r}{281 weeks} \\
\bottomrule
\end{tabular}
\exnote{Weekly strategy returns regressed on an equal-weighted market factor, a
size book and a twelve-week momentum book, all built point-in-time from the same
universe. Newey--West standard errors correct for autocorrelation; first-order
autocorrelation is $+0.063$, so the correction is modest.}
\end{exhibit}

An average beta of zero over seven years does not by itself establish that the
book is neutral when it matters. Exhibit 6 tests it in the period most likely to
break it.

\begin{exhibit}{6}{The 2022 Decline, 61 Weeks}
\begin{tabular}{l d{3.1} d{3.1}}
\toprule
\multicolumn{1}{l}{} & \multicolumn{1}{c}{Market}
 & \multicolumn{1}{c}{Strategy} \\
\midrule
Cumulative return (\%)  & -82.5 & 9.4 \\
Maximum drawdown (\%)   & -84.0 & -7.6 \\
Correlation over window & \multicolumn{2}{c}{0.045} \\
\bottomrule
\end{tabular}

\exnote{Market is the equal-weighted return of the tradeable universe, November 2021 to
December 2022.}
\end{exhibit}

\section{Is It Real?}

Any strategy reported after searching a parameter space needs to answer why the
result is not the best of many tries. We ran four tests, and report the ones that
failed alongside the ones that passed.

\runin{Selection out of sample.}Configuration is chosen on a trailing two-year window and evaluated on the
following six months, rolling forward. Every parameter is reselected in each
window, including the quintile fraction; holding any parameter at its
full-sample best would let the test inherit knowledge of the evaluation period.
Because a weekly strategy also carries hidden exposure to which weekday it
rebalances (Hoffstein, Sibears and Faber 2019), we rebuilt the entire dataset on
each of seven rebalance weekdays and repeated the exercise. We are not aware of
these two procedures being applied together before. Separately each leaves a gap:
a walk-forward result still inherits an arbitrary rebalance weekday, and a
timing-luck study still inherits a fixed configuration. Running them jointly
closes both at once.

Average out-of-sample Sharpe across the seven schedules is 1.87, and all seven are
profitable. Five of the seven clear a $t$-statistic of 3.0, the threshold Harvey,
Liu and Zhu (2016) propose for claimed anomalies. Two do not, and we report them
as failures rather than reporting the best weekday and stopping. The dispersion in
Exhibit 7 is instructive on its own: a study quoting a single schedule could have
reported anything between 1.07 and 2.58 based on a choice with no economic
content.

\begin{exhibit}{7}{Out-of-Sample Result by Rebalance Weekday}
\begin{tabular}{lccccccc}
\toprule
Rebalance day & Sun & Mon & Wed & Sat & Tue & Fri & Thu \\
\midrule
Out-of-sample Sharpe & 2.58 & 2.54 & 1.89 & 1.84 & 1.70 & 1.46 & 1.07 \\
$t$-statistic & 4.82 & 4.75 & 3.53 & 3.44 & 3.17 & 2.73 & 2.01 \\
Clears $t>3.0$ & yes & yes & yes & yes & yes & \textbf{no} & \textbf{no} \\
\bottomrule
\end{tabular}
\exnote{182 out-of-sample weeks per schedule. Average across the seven is 1.87.
The entire dataset is rebuilt from hourly bars on each weekday, so the seven are
genuinely independent schedules rather than shifted views of one. Monday, the
schedule we trade, is second of seven rather than first.}
\end{exhibit}

\runin{Generalisation across assets.}Time-series testing does not show that parameters transfer to contracts they were
not fitted on. We selected configurations on 78 contracts and evaluated on the 34
held out. All eight splits were profitable, retaining 83\% of in-sample
performance once the comparison controls for the fact that a smaller universe
holds fewer positions.

\runin{Correcting for the search itself.}Every configuration evaluated during development was logged to an append-only
record, so the number of trials is counted rather than recalled. Applying the
deflated Sharpe ratio (Bailey and L\'opez de Prado 2014) against the full logged
count of 1,286 trials, the probability that the result reflects skill rather than
selection is 0.95, short of the 0.99 threshold we set in advance. Against the
narrower count of 211 configurations in the final search it is 0.99. We report
both, including the one that fails.

\runin{Costs.}Transaction costs are modelled per contract on trailing dollar volume, at 5 to 25
basis points round trip. The strategy remains profitable at eight times the
modelled level. That is the failure mode most likely to be understated in a
backtest, so it deserves a wide margin.

Exhibit 8 collects the results.

\begin{exhibit}{8}{Robustness Summary}
\begin{tabular}{lcc}
\toprule
Test & Result & Threshold \\
\midrule
Walk-forward, seven rebalance weekdays & Sharpe 1.87 & $>0$ \\
\quad individual weekdays clearing $t>3.0$ & 5 of 7 & --- \\
Held-out contracts, eight splits & 83\% retained & $>50\%$ \\
Deflated Sharpe, 211 trials & 0.99 & 0.99 \\
Deflated Sharpe, 1,286 trials & 0.95 & 0.99 \\
Transaction costs & profitable at $8\times$ & $1\times$ \\
Market loading (Exhibit 5) & $-0.001$ ($t=-0.12$) & $\approx 0$ \\
\bottomrule
\end{tabular}
\exnote{Thresholds are stated in advance rather than chosen after the fact. Two
of the seven rebalance weekdays fall short of $t>3.0$, and the deflated Sharpe
computed against the full logged trial count falls short of 0.99; both are shown
here rather than omitted.}
\end{exhibit}

\section{Does the Live Book Run the Backtest?}

A backtest describes a system that does not exist. The system that eventually
trades is a separate piece of software, usually written later, often by someone
else, and the assumption that the two compute the same positions is almost never
tested. We are not aware of a published strategy paper that verifies it.

We do, with an automated gate that runs before any deployment. Research and the
live path call the same scoring module, and the gate asserts that they produce
identical weights, contract by contract, at every rebalance in the sample. It
also re-tests every factor for look-ahead by recomputing a historical score with
all subsequent data deleted and requiring an exact match, checks that the book is
dollar neutral, and fails if a configuration produces no positions at all. It
exits non-zero on any violation, so a broken build cannot ship.

The gate matters because it has caught real failures. An earlier version of this
work shipped a live book that had drifted to a rank correlation of 0.15 against
research, with position overlap barely above chance. A later look-ahead bug in a
shared ranking helper survived undetected precisely because research and
production agreed perfectly: both called the same incorrect function. That second
failure is the instructive one. Agreement between two systems is not evidence of
correctness if both inherit the same mistake, which is why the gate checks
point-in-time integrity independently rather than only comparing the two paths to
each other.

For a practitioner the implication is direct. The claim that a live system
implements a published backtest is an empirical proposition, not a safe
assumption, and it is cheap to test.

\section{What This Does Not Show}

Three limitations are worth stating plainly, because a reader deciding whether to
act on this should know where the evidence thins out.

The sample covers seven years of a single asset class in what has been, on net, a
period of extraordinary growth and volatility. Seven years is a reasonable sample
in weekly observations but a short one in market regimes. The 2022 decline is the
only severe drawdown in it, and while the strategy came through that episode well,
one episode is one observation.

Capacity is genuinely limited. At roughly \$67 million the strategy is a viable
book for an individual or a small fund and irrelevant to a large one. Readers at
institutions should treat the ranking-frame result as the transferable finding and
the specific strategy as a demonstration of it.

Finally, the result is about construction, not forecasting. We have not shown that
these eleven factors predict returns better than others, or that the effect
persists as more capital pursues it. What we have shown is that given a set of
signals, two decisions taken after signal selection are worth more than the
signals themselves. That is a claim about method, and it should be tested in other
markets that lack a natural peer group before it is treated as general.

\section{Costs, Capacity and Implementation}

The strategy trades 112 contracts weekly with 37\% turnover, which is modest for a
systematic book. All inputs are free and public: prices, volumes and funding rates
come from an exchange bulk archive, positioning statistics from a public endpoint
requiring no credentials. No paid data is used anywhere in the construction.

Capacity is the binding practical constraint. Assuming participation at 2\% of
each contract's trailing dollar volume, the strategy supports roughly \$67
million. That is a meaningful number for an individual manager or a small fund and
an irrelevant one for a large institution. It also implies the effect is unlikely
to be arbitraged away quickly, since the capital required to do so exceeds what
the market will absorb.

One factor requires five-minute data to construct. Practitioners working from
daily bars can implement ten of the eleven; the remaining factor contributes
positively but is not load-bearing.

\section{Conclusion}

Cross-sectional ranking is so standard that its precondition goes unexamined.
That precondition, that the assets being compared are genuinely comparable, is
supplied in equities by accounting and industry structure. Crypto perpetual
futures do not supply it, and ranking them against each other imposes a
comparison the market does not support.

Ranking each contract against its own history removes the requirement. It is not a
refinement of cross-sectional investing but a different premise, and in this
market it is worth 0.44 in Sharpe across eleven independent factors. Holding
factors in separate portfolios rather than blending their signals is worth a
further 0.67. Neither choice requires new data, a proprietary signal or a
forecasting model. Both are decisions about construction, made before any factor
is chosen.

For a practitioner the practical implication is narrower and more useful than the
strategy itself. If you already run a cross-sectional book in this market, the
ranking frame is a change you can test in an afternoon against your existing
signals: replace the cross-sectional percentile with a trailing self-percentile and
rerun. It requires no new data and no new signal research, and in our sample it
was worth more than any individual factor we tested. The construction change,
holding separate books rather than blending scores, is a larger refactor but is
similarly independent of which signals you use.

The broader implication is that in markets without a natural peer group, the
comparison basis is a design decision rather than a convention to inherit. Equity
practice supplies the cross-sectional default, and it supplies it for a good
reason, because equities have the shared structure that makes peer comparison
meaningful. Carrying that default into a market that lacks the structure imports
the machinery without the justification. The same question arises anywhere the
constituents are genuinely heterogeneous: commodities across storability regimes,
credit across capital structures, private assets with no common valuation basis.
The comparison basis deserves the same scrutiny practitioners already give to
signal selection, and in markets like this one it matters more.

\section{Declaration of Submission}

The author has read the publication agreement and is aware of the conditions of
submission. The author declares that no institution requires that this paper be
posted on any public-facing website, including SSRN, ResearchGate or SciHub.

\section{Declaration of Interest}

The author declares no competing financial or personal interests. This research
received no funding. The author has no affiliation with, consulting relationship
with, or financial interest in any cryptocurrency exchange, trading venue or data
provider referenced in this study.

\section{References}

\noindent\hangindent=0.5in Asness, Clifford S., Tobias J. Moskowitz, and
Lasse Heje Pedersen. 2013. ``Value and Momentum Everywhere.'' \emph{The Journal of
Finance} 68 (3): 929--985.

\noindent\hangindent=0.5in Bailey, David H., and Marcos L\'opez de Prado. 2014.
``The Deflated Sharpe Ratio: Correcting for Selection Bias, Backtest Overfitting,
and Non-Normality.'' \emph{The Journal of Portfolio Management} 40 (5): 94--107.

\noindent\hangindent=0.5in Barndorff-Nielsen, Ole E., Silja Kinnebrock, and Neil
Shephard. 2010. ``Measuring Downside Risk: Realised Semivariance.'' In
\emph{Volatility and Time Series Econometrics}, 117--136. Oxford: Oxford
University Press.

\noindent\hangindent=0.5in Chordia, Tarun, Richard Roll, and Avanidhar
Subrahmanyam. 2002. ``Order Imbalance, Liquidity, and Market Returns.''
\emph{Journal of Financial Economics} 65 (1): 111--130.

\noindent\hangindent=0.5in Cong, Lin William, G. Andrew Karolyi, Ke Tang, and
Weiyi Zhao. 2026. ``Crypto Value, Factor Pricing, and Market Segmentation.''
\emph{Management Science}, forthcoming.

\noindent\hangindent=0.5in Fama, Eugene F., and Kenneth R. French. 1993. ``Common
Risk Factors in the Returns on Stocks and Bonds.'' \emph{Journal of Financial
Economics} 33 (1): 3--56.

\noindent\hangindent=0.5in Grinold, Richard C. 1989. ``The Fundamental Law of
Active Management.'' \emph{The Journal of Portfolio Management} 15 (3): 30--37.

\noindent\hangindent=0.5in Harvey, Campbell R., Yan Liu, and Heqing Zhu. 2016.
``$\ldots$ and the Cross-Section of Expected Returns.'' \emph{The Review of
Financial Studies} 29 (1): 5--68.

\noindent\hangindent=0.5in Hoffstein, Corey, Justin Sibears, and Nathan Faber.
2019. ``Rebalance Timing Luck: The Difference Between Hired and Fired.''
\emph{The Journal of Index Investing} 10 (1): 27--36.

\noindent\hangindent=0.5in Kakushadze, Zura. 2016. ``101 Formulaic Alphas.''
\emph{Wilmott} 2016 (84): 72--81.

\noindent\hangindent=0.5in Kyle, Albert S. 1985. ``Continuous Auctions and Insider
Trading.'' \emph{Econometrica} 53 (6): 1315--1335.

\noindent\hangindent=0.5in Lee, Charles M. C., and Mark J. Ready. 1991.
``Inferring Trade Direction from Intraday Data.'' \emph{The Journal of Finance}
46 (2): 733--746.

\noindent\hangindent=0.5in Moskowitz, Tobias J., Yao Hua Ooi, and Lasse Heje
Pedersen. 2012. ``Time Series Momentum.'' \emph{Journal of Financial Economics}
104 (2): 228--250.

\noindent\hangindent=0.5in Newey, Whitney K., and Kenneth D. West. 1987. ``A
Simple, Positive Semi-Definite, Heteroskedasticity and Autocorrelation Consistent
Covariance Matrix.'' \emph{Econometrica} 55 (3): 703--708.

\noindent\hangindent=0.5in Zhou, Wei-Xing. 2008. ``Universal Price Impact
Functions of Individual Trades in an Order-Driven Market.'' \emph{Quantitative
Finance}. Working paper.

\noindent\hangindent=0.5in Patton, Andrew J., and Kevin Sheppard. 2015. ``Good
Volatility, Bad Volatility: Signed Jumps and the Persistence of Volatility.''
\emph{The Review of Economics and Statistics} 97 (3): 683--697.

\end{document}
